% !TeX root = ../main.tex

\begin{abstract}

本論文研究以日常智慧型手機或平板電腦錄製的手部旋轉影片，對帕金森氏症（PD）進行二元分類。所提出的四階段流程包含：(i) 透過 MediaPipe Hands 提取每幀雙手各 21 個三維關鍵點軌跡；(ii) 針對各關節軸軌跡計算頻率、清晰度及四種互補振幅表示（逐週期振幅、滾動峰對峰、滾動均方根、希爾伯特包絡）等運動學特徵，共形成 1,764 維特徵向量；(iii) 以 ANOVA F 值、L1 正則化邏輯斯迴歸或 XGBoost 增益重要性進行特徵選擇；(iv) 以 XGBoost 或特徵式自適應圖卷積風格模型（AGCN-style）作為二元分類器。實驗在臺大醫院兩個臨床世代（2020 及 2025 年）的停藥錄影資料上，分別採用留一法交叉驗證（LOOCV）及 2020 至 2025 年的跨世代遷移協定進行評估。

主要發現如下：第一，左手特徵結合 XGBoost 重要性選擇形成最強的單手基線，且前 10 名選取特徵中，左手佔比在 2020 世代為 25:5，2025 世代為 19:11，與運動儲備假說相符，顯示非慣用手可能呈現更明顯的運動退化。第二，在雙手特徵集中加入關節間配對相關特徵，僅能帶來微小的跨世代 AUROC 提升（$+0.02$--$0.03$），無法穩定改善準確率或 F1 分數。第三，AGCN 在 LOOCV 下 AUROC 約達 $0.98$，但在跨世代遷移中降至 $0.59$，低於最佳手工特徵方法（AUROC $0.68$）。AGCN 的可學習圖拓撲結構高度依賴訓練資料的分布，在小型臨床資料集中難以跨越不同錄製條件與患者族群之間的領域偏移。整體而言，以左手特徵結合 XGBoost 重要性選擇，在 LOOCV 下可達到最高約 82\% 準確率，是目前最穩健的基線方法。

\end{abstract}

\begin{abstract*}

This thesis investigates binary classification of Parkinson's Disease (PD) from short hand-rotation videos recorded with everyday smartphones or tablets. The proposed four-stage pipeline consists of: (i)~3D hand keypoint extraction via MediaPipe Hands, yielding 21 landmark trajectories per hand per frame; (ii)~kinematic feature extraction computing frequency, clarity, and four complementary amplitude representations---cycle-by-cycle, rolling peak-to-peak, rolling RMS, and Hilbert envelope---from each joint-axis trajectory, forming a 1,764-dimensional feature vector; (iii)~feature selection via ANOVA F-value, L1-regularized logistic regression, or XGBoost gain-based importance; and (iv)~binary classification with XGBoost or a feature-based Adaptive Graph Convolutional Network style model (AGCN-style). Five experiments are conducted on two NTUH clinical cohorts (2020 and 2025) of off-medication recordings under Leave-One-Out Cross-Validation (LOOCV) and a 2020-to-2025 cross-cohort transfer protocol.

The principal findings are as follows. First, the left-hand feature set with XGBoost gain-based importance selection forms the strongest single-hand baseline, and the top-10 selected features aggregated across the three feature-selection methods favor the left hand by a 25:5 ratio in the 2020 cohort and 19:11 in the 2025 cohort. This pattern is compatible with the motor reserve hypothesis, but because subject-level handedness labels are not available, it is interpreted as an empirical left-versus-right asymmetry rather than direct proof that the left hand is non-dominant for all subjects. Second, augmenting the bilateral feature set with pairwise joint-correlation features yields only marginal cross-cohort AUROC improvement ($+0.02$--$0.03$) without a stable gain in accuracy or F1-score. Third, AGCN achieves AUROCs of approximately 0.98 under LOOCV but collapses to 0.59 in cross-cohort transfer, falling below the best hand-crafted feature baseline (AUROC 0.68). The sharp LOOCV-to-transfer reversal indicates that the learned graph topology is highly specific to the training distribution and does not generalize across recording protocols or patient populations in small clinical datasets. Overall, combining left-hand features with XGBoost gain-based importance selection achieves up to approximately 82\% accuracy under LOOCV and represents the most robust baseline configuration.

\end{abstract*}
